% !TeX program = xelatex
% 2026 年高教社杯 C 题 —— 摘要
\documentclass[12pt,a4paper]{article}

% ===== 中文支持 =====
\usepackage[UTF8,heading=true]{ctex}

% ===== 页面与行距 =====
\usepackage[a4paper,top=2.5cm,bottom=2.5cm,left=2.8cm,right=2.8cm]{geometry}
\usepackage{setspace}
\setstretch{1.5}

% ===== 数学与符号 =====
\usepackage{amsmath,amssymb}
\usepackage{siunitx}

% ===== 关键词标题格式 =====
\newcommand{\keywordfont}{\bfseries}

\begin{document}

\begin{center}
    {\zihao{3}\heiti 摘\quad 要}\\[1.2ex]
\end{center}

\begin{quotation}
\noindent
微网作为支撑光伏发电就地消纳、缓解并网冲击的重要基础设施，对其购电决策与储能调度问题的研究具有十分重要的意义。本文针对以上问题，根据分时电价、分时段负载与光伏功率及储能参数等数据，并基于季节性朴素预测与 LightGBM 联合预测、混合整数线性规划、滚动优化与多阶段调整等研究方法，对不确定供需下购电与储能的耦合调度规律进行详细分析，建立覆盖预测、计划、修订与评估四个环节的购电决策模型。

\medskip
\noindent\textbf{针对问题一：}在电价与负载固定、各时段光伏发电功率已知的条件下，建立了以全天购电费用最小化为目标的 MILP 模型。模型以 144 个 10 分钟时段为决策周期，明确购电量、电池贮存电量、充放电 $0$--$1$ 变量等决策变量，精确刻画储能设备物理特性，最后采用 HiGHS 求解器求解。最终得到全天最优购电费用为 $35\,126.85$ 元，购电量 $59\,482.70$~kWh，相较于无储能方案节省 $26.90\%$。

\medskip
\noindent\textbf{针对问题二：}在负载与光伏实时波动、供需缺口须按 $5$ 倍电价应急购电的条件下，本文先依据周内与日间规律构造季节性朴素预测，然后引入 $15$ 个负载特征与 $12$ 个光伏特征并用 LightGBM 修正朴素预测残差，并通过历史残差分位数保护预测不确定性，验证了负载与光伏显著的日内周期性、负载更强的周周期性以及光伏相邻日之间的连续性。在此基础上，建立了以计划购电费用最小为目标、受电力供需平衡、每日始末 SOC 一致、SOC 上下限及充电功率上限等约束的混合整数线性规划模型：
\begin{equation*}
\boxed{\;\min\; C_{\mathrm{plan},k}=\sum_{j=0}^{143} p_j\,q_{k,j}\;}
\end{equation*}
其中 $p_j$ 为第 $j$ 时段分时电价，$q_{k,j}$ 为第 $k$ 天第 $j$ 时段计划购电量。计算得优化后全年购电费用 $xx$ 元，相比优化前 $xx$ 元降低 $xx\%$，应急购电费用占比 $xx\%$。

\medskip
\noindent\textbf{针对问题三：}本文在问题二日前优化计划基础上，构建了考虑计划调整违约成本、预测误差保护和储能 SOC 实时反馈的日内多阶段滚动调度模型：按照题目规定的 $00{:}00$、$06{:}00$、$12{:}00$ 和 $18{:}00$ 四个预报更新节点，采用分段线性插值将光伏预测转换为 10 分钟尺度数据，并利用已观测负荷误差对后续负荷预测进行前向自适应修正，根据修正后净需求误差的 $75\%$ 分位数配置安全保护量；每次新预测到达后，以当前真实 SOC 为初始状态重新优化剩余时段，并比较原计划与新计划的普通购电费、调整违约费和预期应急购电费，仅在综合成本更优时执行调整。在 2025 年评价期和固定参数配置下，模型全年总费用降至 $13\,330\,377.96$ 元，较日前基准方案节省 $23\,444.08$ 元，降幅约 $0.1756\%$；应急购电电量由 $36\,526.87$~kWh 降至 $20\,406.68$~kWh，减少 $44.13\%$，应急购电费用由 $18.62$ 万元降至 $10.50$ 万元，减少 $43.60\%$。结果表明，该滚动调度机制能够以一定计划调整成本为代价降低应急购电需求，并实现购电综合费用下降。

\medskip
\noindent\textbf{针对问题四：}针对外部电网电价逐日逐段波动的特点，本文将波动电价引入问题二和问题三的调度模型，对两类调度方案进行重新求解。结果表明，问题二重新求解后的全年运行费用为 $14\,347\,907.89$ 元，问题三结合日内信息更新进行滚动重优化后的全年运行费用为 $13\,972\,808.91$ 元，后者较前者减少 $375\,098.98$ 元，降幅为 $2.6143\%$。结果说明，在波动电价环境下重新优化购电时机与储能充放电策略，能够降低微网运行费用；同时，调度结果仍需结合调整费用和应急供电风险进行综合评价。

\end{quotation}

\vspace{2em}
\noindent\keywordfont{关键词：} 微网购电决策；储能调度；混合整数线性规划；LightGBM 预测；日内滚动优化

\end{document}
